\chapter{Methodology}
\label{chap: methodology}

[INTRODUCTION AND SECTION OUTLINE]\\

\section{The Supply Chain Model as a Cooperative Partially Observable Markov Game}
\label{sec:supply-chain-pomg}
In order to consistently apply \gls{acr:marl} methods, we model the distribution network of \autoref{chap:problem-formulation} as a finite-horizon, discounted \gls{acr:pomg}. In this section, we thus specify the central \gls{acr:pomg} components in the context of the distribution network and outline our design of the observation functions $O^i$, the action space $\mathcal{A}$, and the rewards $R^i$.\\

\subsection{State and Observation Design}
\label{subsec:observation-design}
We treat the warehouses as the agents of the game such that $\mathcal{I}=\mathcal{W}$ and agent $i$ is warehouse $i$. The game therefore has $N=W$ agents that need to make decisions on their replenishment orders by interacting with the environment, which is the distribution network of \autoref{chap:problem-formulation}. In order to apply a deep learning architecture to this setting, we must formally specify a collection of features that agents observe and act on. Broadly, this can be approached in two ways. 

One method favors a minimal feature set to obtain a minimum viable observation $o^{i,\min}_t$ that contains only the most essential information needed to describe each agent's local situation. The idea behind this approach is that a sufficiently deep \gls{acr:nn} could itself learn an effective policy by abstracting complex relationships from these basic features such that explicit feature engineering becomes unnecessary [\textbf{CITE: paper that shows this}].
For instance, a minimal feature set for our setting would expose only the current inventory state of the agent, which is fully characterized by its on-hand inventory $I_{ikt}$ and its in-transit orders by arrival slot $\mathbf{I}^{\mathrm{tr}}_{ikt} = \left( q_{ik,\,t-L_{ik}+1},\ \dots,\ q_{ikt} \right)$, so that $ o_{it}^{min} = [ (I_{ikt})_{k \in \mathcal{K}},\ (\mathbf{I}^{\mathrm{tr}}_{ikt})_{k \in \mathcal{K}} ]$ for all agents. In fact, for a single warehouse under stationary demand, lost sales, and a fixed lead time the inventory state is a sufficient statistic for the optimal ordering policy [\textbf{CITE: inventory position as a sufficient statistic, {karlin1958lostsales}}].

However, since we consider the multi-agent case in which the agents are coupled through transshipment, a single agent's local state is no longer a sufficient statistic and we can thus not rely on the single-installation result of [\textbf{CITE:{karlin1958lostsales}}]. Additionally, both prior work and our own experiments indicate that incorporating a broader set of features significantly improves convergence speed, final policy performance, and model interpretability [\textbf{CITE: benefits of feature engineering in (multi-agent) RL for inventory control, {glanois2024survey}}]. As part of our feature engineering process, we have therefore designed a more refined set of features that incorporates our prior understanding of the environment into the model and summarizes recent history better compared to a single-period snapshot. Specifically, we categorize our features into four categories that follow the structure of the environment: \textit{inventory features}, \textit{demand features}, \textit{fulfillment features}, and \textit{identity features}. While inventory features capture the stock a warehouse currently holds or has in transit, demand features summarize the home-region demand a warehouse observes. Fulfillment features record how much of incoming demand a warehouse served and where it shipped and identity features identify which warehouse an observation belongs to. A comprehensive list of the features used in our implementation is provided in \autoref{tab:obs-features}. 

\begin{table}[tbp]
\centering
\small
\setlength{\tabcolsep}{6pt}
\renewcommand{\arraystretch}{1.3}
\caption{Local observation features grouped by category, with the symbol and definition
for each. All quantities are for warehouse $i$, \gls{acr:sku} $k$, and period $t$ unless
noted. For home-region quantities the region index is set to $j=i$.}
\label{tab:obs-features}
\begin{tabular}{@{}l l >{\raggedright\arraybackslash}m{0.66\linewidth}@{}}
\toprule
\textbf{Category} & \textbf{Symbol} & \textbf{Definition}\\
\midrule
\multirow{5}{*}{\textbf{Inventory}}
 & $I_{ikt}$ & On-hand inventory\\\addlinespace[1pt]
 & $I^{\mathrm{agg}}_{it}$ & Total on-hand inventory across all \glspl{acr:sku}, $\sum_{k} I_{ikt}$\\\addlinespace[1pt]
 & $I^{\mathrm{net}}_{ikt}$ & Expected net inventory after the lead time, $I_{ikt}+I^{\mathrm{tr}}_{ikt}-L_{ik}\hat{D}_{ikt}$\\\addlinespace[1pt]
 & $\mathrm{DoS}_{ikt}$ & Days of supply based on recent average demand, $I_{ikt}/\bar{D}_{ikt}$\\\addlinespace[1pt]
 & $\mathbf{I}^{\mathrm{tr}}_{ikt}$ & Outstanding orders by arrival slot, summing to the pipeline $I^{\mathrm{tr}}_{ikt}$\\\addlinespace[1pt]
\midrule
\multirow{5}{*}{\textbf{Demand}}
 & $D_{ikt}$ & Incoming home-region demand\\\addlinespace[1pt]
 & $\bar{D}_{ikt}$ & Mean home-region demand over the last $5$ periods\\\addlinespace[1pt]
 & $\hat{D}_{ikt}$ & Exponentially smoothed home-region demand ($\alpha=0.3$)\\\addlinespace[1pt]
 & $\sigma^{D}_{ikt}$ & Standard deviation of home-region demand over the last $5$ periods\\\addlinespace[1pt]
 & $\mathbf{D}_{ikt}$ & Incoming home-region demand in each of the last $5$ periods\\\addlinespace[1pt]
\midrule
\multirow{2}{*}{\textbf{Fulfillment}}
 & $S^{\mathrm{home}}_{ikt}$ & Allocated home-region shipments\\\addlinespace[1pt]
 & $S^{\mathrm{away}}_{ikt}$ & Allocated foreign-region shipments\\\addlinespace[1pt]
\midrule
\textbf{Identity}
 & $\mathbf{e}_{i}$ & One-hot vector of length $W$ identifying warehouse $i$\\\addlinespace[1pt]
\bottomrule
\end{tabular}
\end{table}

We approached the feature design with three central objectives in mind. First, we aimed to improve convergence speed and stability by designing features that simplify action selection. Take, for instance, the pipeline vector $\mathbf{I}^{\mathrm{tr}}_{ikt}$. By tracking the arrival slots of units it has already ordered but not yet received, the agent can optimize its behavior through ordering only the remaining shortfall rather than reordering stock that is already in transit. Second, we intended to provide the agents with deeper insights into the environment by combining, and in some places replacing, multiple raw inputs into a single more meaningful feature. One example is the days of supply $\mathrm{DoS}_{ikt}$, which expresses an agent's stock as the days of average recent demand it can cover. By combining inventory and average demand, we expected this feature to help the agent judge whether it is over- or understocked more effectively at each decision step. Similarly, the expected net inventory $I_{ikt}^{\mathrm{net}}$ gives the stock that the agent can expect to hold once all outstanding orders have arrived and lead-time demand has been served and thus might function as an early indicator of potential shortfalls. Finally, these engineered features allow us to reconcile single-observation networks with the history-dependent policy class defined in [\textbf{section POMG}]. By embedding past demand dynamics directly into the current observation through features such as $\mathbf{D}_{ikt}$ or $\bar{D}_{ikt}$, a feedforward network reading only $o^i_t$ can successfully approximate a policy conditioned on the history $\mathbf{h}_t^i$.

However, it is important to note that we developed these features based on our empirical intuition rather than a theoretical framework. While they are designed to enhance learning and decision-making, there is no guarantee that each feature will improve performance. Indeed, some may duplicate information that another feature
already captures and others may enhance training efficiency at the expense of generalization and robustness. For this reason, we do not fix the feature set in advance and instead make feature selection a central part of our tuning process (see \textbf{autoref{sec:hyperparameter-tuning}}). The only two features we keep active in every configuration are the on-hand inventory $I_{ikt}$ and pipeline vector $\mathbf{I}^{\mathrm{tr}}_{ikt}$ since we consider them the basis of any ordering decision.

Formally, the environment evolves on the full physical state of \autoref{sec:replenishment},
which stacks the on-hand inventory and the outstanding orders by arrival slot of every
Markov property in that the
probability of the next state and reward is conditionally independent of the past
states and joint actions,
warehouse and \gls{acr:sku},
\begin{equation}
\label{eq:state}
  s_t = \big[\, (I_{ikt})_{i \in \mathcal{W},\, k \in \mathcal{K}},\ (\mathbf{I}^{\mathrm{tr}}_{ikt})_{i \in \mathcal{W},\, k \in \mathcal{K}} \,\big].
\end{equation}
This state is Markov for the joint dynamics, since demand is independent across periods and the transition depends on the past only through the current inventory and pipeline. No agent observes $s_t$ in full and each instead acts on a local observation assembled from the features of \autoref{tab:obs-features}. we partition the features of \autoref{tab:obs-features} into the default inputs $I_{ikt}$, $\mathbf{I}^{tr}_{ikt}$, $\mathbf{e}_i$ and a catalogue of tuneable blocks $\{\psi_1, \dots, \psi_M\}$, where each block $\psi_m$ returns a real vector from agent $i$'s local history $\mathbf{h}^i_t$ and contains either the latest entry or a trailing window of recent observations from $\mathbf{h}^i_t$, depending on the feature. This catalogue serves as the basis for our tuning process, which selects features by choosing a specific configuration that activates a subset of blocks $\mathcal{B} \subseteq \{1, \dots, M\}$. The local observation of agent $i$ is then the concatenation of the default parameters and the active blocks in a fixed order,
\begin{equation}
\label{eq:observation}
    o^i_t = \big[\, \mathbf{e}_i,\ (I_{ikt})_{k \in \mathcal{K}},\ (\mathbf{I}^{\mathrm{tr}}_{ikt})_{k \in \mathcal{K}},\ \big( \psi_m(\mathbf{h}^i_{t-1}) \big)_{m \in \mathcal{B}} \,\big] \in \mathbb{R}^{d},
\end{equation}
where $d$ is the summed width of the default and active blocks. Under this convention, each agent $i$ has a single local observation $o^i_t$ that is a deterministic function of its current inventory and pipeline and of its past local history $\mathbf{h}^i_{t-1}$. Since all $W$ local observations share identical features, concatenating them yields the global observation $z_t = (o^1_t, \dots, o^W_t)$, which can directly be used by a centralized critic.\\


\subsection{Action design}
\label{subsec:action-design}
At each time step $t$, each agent $i$ must determine the replenishment quantity $q_{ikt} \ge 0$ for each \gls{acr:sku} $k$ based on its most recent local observation $o_{t-1}^{i}$. While there are various possible formulations of the action space, a naive approach would have the policy directly output $q_{ikt}$ on a bounded scale. However, we avoid this direct design for two reasons. First, raw quantities can render policy initialization and early exploration less effective. Since raw-quantity actions lack a physical anchor to demand, a randomly initialized policy would output arbitrary quantities at startup that are decoupled from actual environment requirements. The agent must then spend early training episodes on calibrating actions to a sensible range, which leads to inefficient exploration and reduces sample efficiency. Second, raw quantities can cause parameter sharing to suffer in heterogeneous settings. Specifically, in settings where demand profiles differ across warehouses or \glspl{acr:sku}, a shared policy must learn agent-specific scale factors across an output range that is wide enough to cover all demand levels simultaneously. This not only slows training convergence but also increases sensitivity to prediction errors as minor absolute errors in policy outputs can translate into substantial misallocations.

To address these issues, we adopt the $X+Y$ replenishment policy of [\textbf{CITE: Kimbrough et al. (2002), Preil and Krapp (2022)}]. According to this rule, the order quantity $q_{ikt}$ for each \gls{acr:sku} $k$ and agent $i$ is determined by the two values $X$ and $Y$, where $X$ is the most recently observed home-region demand $D_{ik,t-1}$ and $Y$ is some variable adjustment $y_{ikt}$ to $D_{ik,t-1}$. Following our notation and adopting the $X+Y$ rule, the order quantities are thus defined as 
\begin{equation}
\label{eq:xy-action}
  q_{ikt} = \max\!\big\{0,\ D_{ik,t-1} + y_{ikt}\big\},
\end{equation}
where the $\operatorname{max}$ operator ensures that orders are non-negative and we set $D_{ik,-1} := 0$ so that the first-period order reduces to $q_{ik,0} = \max\{0,\ y_{ik,0}\}$. Since the deviations $y_{ikt}$ have to be determined for each agent $i$, \gls{acr:sku} $k$, and time $t$ separately, it follows that they replace the order quantities $q_{ikt}$ as the primary decision variables for each agent $i$ such that the final replenishments are only derived implicitly through \autoref{eq:xy-action}.

Adopting the $X+Y$ rule directly resolves the previously discussed limitations of raw-quantity actions. Because the policy now controls only the deviation $y_{ikt}$, a randomly initialized policy explores relative adjustments to demand rather than the order quantity itself. This keeps proposed orders within a bounded range near observed demand and thus skips the early training phase needed to learn demand scales. Expressing the action as a deviation from observed demand likewise removes the scale-dependence issue since the action is centered on observed demand and thus carries the same meaning across differing warehouses or \glspl{acr:sku} irrespective of absolute demand. Beyond addressing these two issues, the relative structure of the $X+Y$ rule offers an additional and arguably more important benefit. By restricting each deviation $y_{ikt}$ to the fixed range $\{-\Delta_k,\dots,\Delta_k\}$, we can effectively reduce the action space from all possible quantity vectors to a smaller set of demand-relative adjustments that the policy can search far more efficiently [\textbf{CITE: action-space reduction benefit of the $X+Y$ rule, preil and krapp mcts}].

However, choosing $\Delta_k$ requires balancing two opposing needs. While $\Delta_k$ must be sufficiently large to provide the agent with ordering flexibility, it must also remain narrow enough to preserve the reduced action space that motivated the rule. Moreover, since the appropriate size of $\Delta_k$ depends on the demand scale that can differ across \glspl{acr:sku} and environments, no single fixed value for $\Delta_k$ can balance both needs in all cases. To address these considerations, we set $\Delta_k$ using the following rule, 
\begin{equation}
\label{eq:delta-rule}
  \Delta_k = \max\!\big\{2,\ \lceil 1.5 \cdot \max_{i \in {\mathcal{W}}} \mathbb{E}[D_{ikt}] \rceil\big\},
\end{equation}
where $\mathbb{E}[D_{ikt}]$ is the expected per-step home-region demand for \gls{acr:sku} $k$ at warehouse $i$. To allow the agent to withhold ordering during periods of overstock, the lower bound of the adjustment range must allow for a net order of zero under typical demand. Therefore, our rule enforces the condition $\Delta_k \ge \mathbb{E}[D_{ikt}]$ to ensure that the minimum possible order under average demand, $\max\!\big\{0,\ \mathbb{E}[D_{ikt}] - \Delta_k\big\}$, can drop to zero. Similarly, the upper bound must give the agent the necessary upward flexibility to accumulate stock for future periods, while remaining restrictive enough to preserve a limited action space. To achieve this, we scale the expected demand by a factor of $1.5$ as this multiplier proved effective in early experiments. Note that this scaling factor does not violate the requirement for zero-orders under average demand. The remaining operators in \autoref{eq:delta-rule} handle network and scale considerations. Taking the maximum over $\mathcal{W}$ ensures that a single shared $\Delta_k$ remains sufficient for the highest-demand warehouse and applying the outer $\max\{2,\cdot\}$ holds the action range at a minimum of $\pm2$ so that \glspl{acr:sku} with very small $\mathbb{E}[D_{ikt}]$ do not collapse to a near-trivial action space.
Since the learned decision variables are the adjustments $y_{ikt}\in\{-\Delta_k,\dots,\Delta_k\}$ instead of the absolute order quantities $q_{ikt}$, the local and joint action spaces are defined as
\begin{equation}
\label{eq:action-spaces}
  \mathcal{A}^i = \prod_{k \in \mathcal{K}} \{-\Delta_k,\ \dots,\ \Delta_k\}, \qquad \mathcal{A} = \prod_{i \in \mathcal{W}} \mathcal{A}^i,
\end{equation}
where $\Delta_k$ is shared across warehouses.\\


\subsection{Reward Design}
\label{subsec:reward-design}
Once all agents have observed $o^i_{t-1}$ and selected $y_{ikt}$ to form the joint action $\mathbf{a}_t \in \mathcal{A}$ at the start of each period $t$, the environment transitions from the physical state $s_t$ to $s_{t+1}$ with probability $P(s_{t+1} \mid s_t, \mathbf{a}_t)$ according to the \gls{acr:pomg} process. This transition is stochastic since demand is assumed to be generated by the Poisson process of \autoref{sec:demand-model} and is realized only after orders have been placed. Therefore, the agents cannot reliably anticipate the incoming demand such that they are uncertain about the subsequent state that the environment will move to. 
After transitioning to state $s_{t+1}$, the environment generates a reward $r_t^i$ according to the reward function $R^i(s_t,\mathbf{a}_t,s_{t+1})$. We adopt an agent-scoped reward in which each warehouse receives the negative of its own per-period cost,
\begin{equation}
\label{eq:reward}
  r^i_t  = -C_{it}.
\end{equation}
Note that the negative sign aligns reward maximization with cost minimization such that an agent earns a higher reward whenever it lowers its own cost. Because each cost component is attributable to a single warehouse, this per-agent reward is well defined and aggregates naturally across the network to the negative system cost, $-C_t$, and, when accumulated over time, to the negative $\mathrm{\gls{acr:tscc}}$. Therefore, the agents are cooperative and jointly aim to minimize the expected $\mathrm{\gls{acr:tscc}}$, even though each warehouse receives only its own cost. 

The agent-based reward design of \autoref{eq:reward} is only one of two natural ways to distribute the system cost across agents. An alternative, team-based formulation would return the full system cost $C_t$ as a common reward to each agent to achieve equal accountability for the network outcome. However, we restrict our main study to the agent scope because it not only aligns with the real-world data that our distribution model is based on but also holds each warehouse accountable for its own decisions. In turn, this keeps per-agent credit assignment clean under heterogeneity, which is a precondition for attributing potential failure modes to individual warehouses. We retain the team scope as a secondary variant and evaluate its effect in [\textbf{REF: reward-scope robustness experiment}].\\


\subsection{Initial Settings}
\label{subsec:initial-settings}
The observation, action, and reward designs above describe how an episode evolves once it is running but leave open the state from which it starts. This start state propagates through the early transitions and shapes the warehouses' behavior before they settle into their operating regime. We therefore complete the specification of the game by stating how each episode is initialized.

The central initialization decision concerns the on-hand inventory $I_{ik,0}$ that each warehouse $i$ holds for every \gls{acr:sku} $k$ at the start of an episode. Because this inventory governs the first periods of every episode and the warehouses across our settings face demand of very different magnitudes, no single fixed value is appropriate for all of them. A value suited to a high-demand warehouse would leave a low-demand warehouse heavily overstocked at the start, whereas a value suited to a low-demand warehouse would leave a high-demand warehouse short. A shared constant would therefore advantage some settings and handicap others. We thus set the initial inventory through a demand-dependent rule rather than a fixed quantity. Specifically, we initialize each warehouse at the analytical steady-state base stock,
\begin{equation}
\label{eq:initial-inventory}
  I_{ik,0} = \mathbb{E}[D_{ik}]\,L_{ik} + z\,\sqrt{\mathbb{E}[D_{ik}]\,L_{ik}},
\end{equation}
where $\mathbb{E}[D_{ik}]$ is the expected per-period home-region demand of \autoref{eq:expected-demand} and $z$ is a safety factor. The first term places the warehouse at its expected lead-time demand and the second term adds a safety margin scaled by $\sqrt{\mathbb{E}[D_{ik}]\,L_{ik}}$, which is the standard deviation of lead-time demand under a Poisson approximation. We set $z = 3.09$, which is deliberately large. The purpose of this value is operational rather than statistical. A small initial inventory would force the warehouses to spend the first lead times stocking up toward their operating range. This cold-start transient would then dominate the reported return rather than reflecting the behavior of the learned policy. A large $z$ instead starts each episode close to the steady-state operating point so that the reported return reflects the policy in operation. Because \autoref{eq:initial-inventory} adapts to each warehouse and \gls{acr:sku} through its expected demand, it lets us choose initial inventories consistently across the heterogeneous settings introduced later rather than retuning a constant for each one.

This Poisson approximation understates the true variability of the order-level demand, which is over-dispersed according to \autoref{eq:demand-variance}. Since the rule only needs to start each warehouse near its operating point rather than realize an exact service level, the approximation is immaterial for initialization and the conservative $z$ keeps the initial inventory well above the expected lead-time demand in any case. We return to the consequences of this over-dispersion when we introduce the environments in [\textbf{REF: environment design}].

Two further quantities complete the description of the first period. First, the pipeline starts empty such that no replenishment orders are in transit and the pipeline vector satisfies $\mathbf{I}^{\mathrm{tr}}_{ik,0} = \mathbf{0}$. With a fixed lead time the first arrivals reach a warehouse only at period $L_{ik}$ such that the demand during the first $L_{ik}$ periods is served from the initial inventory alone. This is exactly the lead-time demand that \autoref{eq:initial-inventory} is built to cover. Second, the $X+Y$ rule requires a demand anchor for its first decision even though no demand has been observed before the episode begins. As stated in \autoref{subsec:action-design}, we therefore set $D_{ik,-1} := 0$ so that the first order reduces to $q_{ik,0} = \max\{0,\,y_{ik,0}\}$ of \autoref{eq:xy-action}.\\


\section{Baselines and Algorithms}
\label{subsec:baselines-algorithms}


\section{Tuning and Evaluation Protocol}
\label{subsec:tuning-evaluation}




